% Propuesta separada; no modifica el bloque común del artículo I.
% Procedencia: artículo II, revisión 05, registro_incidencias.tex y
% registro_imagen_integral.tex; aclaraciones de índices de revisiones 06/07.
% Inserción prevista inmediatamente antes de subsec:k-terminal.
\subsection{Evaluation of incidences and determination of the transverse coordinates}
\label{iii:r25:incidencias}

The dodecaphasic register preserves a history of visits and oriented
traces. Its transverse coordinates are obtained through a composition
of operations: arithmetic evaluation of the visits, counting by windows,
decimal normalization, and cyclic differences.
Reconstructing the register from those differences is another operation,
whose existence and uniqueness will be proved below.
This separation makes it possible to specify both the data used by the
extractor and the information it transmits to the harmonic readout.

Let $\mathscr L$ be an incidence register produced by the selection,
transport, and updating of an APP--TRIT--TPK history.
Each visit preserves its route, APP position, sheet, time,
multiplicity, and boundary class; each trace preserves its members,
orientation, multiplicity, and length in the unit declared by its
readout. The windows are
\[
 E_m=\{9(m-1)+1,\ldots,9m\},\qquad 1\leq m\leq12.
\]
Evaluation of the sum or product sheet determines the integer
$n(v)$ of each visit and its positive division
\[
 n(v)=r(v)+9q(v),\qquad 1\leq r(v)\leq9.
\]
The residue is calculated from that evaluation. The boundary class of
a visit with residue $9$ is preserved by
$\epsilon_\partial(v)=1$ on the corona and $-1$ in the interior.

The incidence presentation of the register uses the three counts
\begin{align}
 A_m&=\sum_{v\in E_m}w(v)
       \mathbf1_{\{1,3,5,7\}}(r(v)),\label{iii:r25:recuento-a}\\
 C_m&=\sum_{j\geq0}
       \bigl(\nu^-_{120,m}(j)-\nu^+_{120,m}(j)\bigr),
       \label{iii:r25:recuento-c}\\
 V_m&=\sum_{v\in E_m}w(v)
       \mathbf1_{\{9\}}(r(v))\epsilon_\partial(v).
       \label{iii:r25:recuento-v}
\end{align}
Here $w(v)$ is the incidence multiplicity and
$\nu^\pm_{120,m}(j)$ counts the oriented traces of the indicated
channel with reduced length $120(1+3j)$, according to the enumeration
of the register.
The collection of traces and its multiplicity are data of the evaluation,
together with their corresponding production rule. The integer sum requires
finite support or a construction that determines its meaning.
The reduced length of a trace and the number of times in the
window are differently typed quantities; the readout unit specifies
their relationship. The symbols $A_m,C_m,V_m$ here denote
integer counts, distinct from the angular pair and the constitutive operator.

Let $[x]_{10}\in\{0,\ldots,9\}$ be the Euclidean representative.
Normalization also preserves the quotient in
\[
 x=[x]_{10}+10\lfloor x/10\rfloor.
\]
The block is formed as
\begin{equation}
 z_m=100[A_m]_{10}+10[C_m]_{10}+[V_m]_{10},
 \qquad 0\leq z_m\leq999.
 \label{iii:r25:bloque}
\end{equation}
With cyclic indices modulo $12$, let
\[
 (D_a z)_m=z_m-z_{m+a},\qquad a\in\{3,4\},
 \qquad Q(z)=\sum_{m=1}^{12}z_m.
\]
The incidence composition is expressed by
\begin{equation}
 \mathcal E_{\rm cnt}(\mathscr L)
 =\mathcal D(z):=(D_3z,D_4z,Q(z))\in\mathbb Z^{25}.
 \label{iii:r25:composicion}
\end{equation}
Its arguments are the visits and traces with their counting rules.
The evaluation does not receive $K$, $U$, $\alpha$, or the expected
transverse coordinates. To identify it with
$\mathcal E_{108}^{90,120}$, the composition with the actual
collection and incidences of the terminal state is also preserved. The
invertibility of $\mathcal D$ determines a subsequent reconstruction;
by itself, it does not select that collection of incidences.

\begin{proposition}[Residual information determined by the extractor]
\label{iii:r25:residuos}
Let $N,N'\in\mathbb Z^{12\times3}$ be two count registers,
ordered by $(A,C,V)$. For the operations
\eqref{iii:r25:bloque}--\eqref{iii:r25:composicion},
\[
 \mathcal E_{\rm cnt}(N)=\mathcal E_{\rm cnt}(N')
 \quad\Longleftrightarrow\quad N-N'\in10\mathbb Z^{36}.
\]
The $36$ sectoral residues determine exactly the output of
$25$ coordinates. The quotients and provenance constitute
additional information in the register.
\end{proposition}
\begin{proof}
Equality modulo $10$ produces the same blocks and the same
transverse coordinates. Conversely, equality of the
coordinates gives $D_3(z-z')=D_4(z-z')=0$. Steps $3$ and $4$
generate the cycle of length $12$, so $z-z'$ is constant.
Equality of $Q$ makes that constant zero. The expansion of each integer
from $0$ to $999$ into three decimal digits is unique; hence
the three residues of each window are recovered.
\end{proof}

The proposition is valid on the entire indicated integer domain. It
does not postulate that the $36$ classes are independent on the subset
of admissible registers produced by TPK.

\subsection{Integral image and inversion of the transverse coordinates}
\label{iii:r25:imagen-seccion}

To express the inversion, the blocks are reindexed by
$i\in\mathbb Z/12\mathbb Z$, from $0$ to $11$. Let
$(Sz)_i=z_{i+1}$, $D_3=I-S^3$, and $D_4=I-S^4$.
Given an integer triple $(b,c,q)$, define
\begin{equation}
 w_i=c_i-b_{i+1},\qquad
 P_0=0,\quad P_i=\sum_{j=0}^{i-1}w_j\ (1\leq i\leq12),
 \qquad T(w)=\sum_{i=0}^{11}P_i.
 \label{iii:r25:incrementos}
\end{equation}

\begin{theorem}[Characterization of the integral image]
\label{iii:r25:imagen}
A triple $(b,c,q)\in\mathbb Z^{25}$ belongs to
$\operatorname{im}\mathcal D$ if and only if
\begin{align}
 b_i&=w_i+w_{i+1}+w_{i+2},\qquad 0\leq i<12,
       \label{iii:r25:relaciones}\\
 \sum_{i=0}^{11}w_i&=0,\label{iii:r25:cierre}\\
 q+T(w)&\equiv0\pmod {12}.\label{iii:r25:congruencia}
\end{align}
Its unique integer preimage is
\begin{equation}
 z_i=\frac{q+T(w)}{12}-P_i,\qquad0\leq i<12.
 \label{iii:r25:inversa}
\end{equation}
The twelve relations \eqref{iii:r25:relaciones} and relation
\eqref{iii:r25:cierre} are linearly independent over
$\mathbb Q$.
\end{theorem}
\begin{proof}
For an output of $\mathcal D$,
\[
 c_i-b_{i+1}
 =(z_i-z_{i+4})-(z_{i+1}-z_{i+4})=z_i-z_{i+1}.
\]
Three consecutive increments sum to $b_i$, and their total
cyclic sum is zero. Moreover, $P_i=z_0-z_i$, whence
$q=12z_0-T(w)$. This gives the conditions and the inverse formula.
Conversely, the congruence makes the coordinates of
\eqref{iii:r25:inversa} integers, and closure allows them to be
extended cyclically. Their adjacent differences are $w_i$, so $D_3z=b$
and
\[
 (D_4z)_i=w_i+w_{i+1}+w_{i+2}+w_{i+3}
 =w_i+b_{i+1}=c_i.
\]
Summation gives $Q(z)=q$. The formula also determines uniqueness.
Finally, the invertible integer change
$(b,c)\mapsto(b,w=c-Sb)$ transforms the first twelve relations
into $b=(I+S+S^2)w$: each solves for a different coordinate of
$b$. Closure $\sum_iw_i=0$ adds an independent relation.
\end{proof}

The rational image has dimension $12$. The integral condition adds
the congruence of order $12$ in \eqref{iii:r25:congruencia}.
The eleven increments $w_0,\ldots,w_{10}$ and the charge $q$ are
sufficient coordinates, with $w_{11}=-\sum_{i=0}^{10}w_i$.
The $25$ coordinates preserve those $12$ degrees of freedom
through $13$ linear relations. If the triple also comes from
\eqref{iii:r25:bloque}, the inverse belongs to
$\{0,\ldots,999\}^{12}$ and recovers its three residues per window.

\subsubsection{Agreement with the harmonic transform}

The transform $\mathcal H_{12}$ acts through $H_4$ on the
orbits $(r,r+3,r+6,r+9)$, $r=0,1,2$, where
\[
 H_4=\begin{pmatrix}
 1&1&1&1\\1&1&-1&-1\\1&-1&1&-1\\1&-1&-1&1
 \end{pmatrix},\qquad H_4^2=4I_4.
\]
For a compatible triple, let
\[
 B_r=\sum_{j=0}^3c_{r+3j},\qquad
 a_0=\frac{q+2B_0+B_1}{3},\quad
 a_1=a_0-B_0,\quad a_2=a_1-B_1.
\]
The harmonic readout has blocks
\[
 (\Pi_H(b,c,q))^{(r)}
 =(a_r,b_{r+3}-b_{r+9},b_r+b_{r+6},b_r-b_{r+6}).
\]

\begin{proposition}[Agreement of reconstructions]
\label{iii:r25:hadamard}
On $\operatorname{im}\mathcal D$,
\[
 \Pi_H\mathcal D=\mathcal H_{12},\qquad
 \mathcal D^{-1}=\tfrac14\mathcal H_{12}\Pi_H.
\]
The image of $\Pi_H$ restricted to that domain is
\[
 \mathcal L_H=
 \{u\in\mathbb Z^{12}:\mathcal H_{12}u\equiv0\pmod4\}.
\]
\end{proposition}
\begin{proof}
For $z=\mathcal D^{-1}(b,c,q)$, let
$h_r=\sum_{j=0}^3z_{r+3j}$. Shifting by four positions
takes each orbit to the next; therefore $B_r=h_r-h_{r+1}$.
The identity $q=h_0+h_1+h_2$ gives $a_r=h_r$.
Within an orbit, write $(x_0,x_1,x_2,x_3)$ for its coordinates
and $d_j=x_j-x_{j+1}$. Then
\[
 \begin{aligned}
 d_1-d_3&=x_0+x_1-x_2-x_3,\\
 d_0+d_2&=x_0-x_1+x_2-x_3,\\
 d_0-d_2&=x_0-x_1-x_2+x_3.
 \end{aligned}
\]
These are the remaining three rows of $H_4x$. Inversion and the
integral characterization follow from $\mathcal H_{12}^2=4I$.
\end{proof}

The harmonic congruence condition controls the output of $\Pi_H$;
membership of the channels in $\operatorname{im}\mathcal D$
also requires \eqref{iii:r25:relaciones}--\eqref{iii:r25:congruencia}.
An exact negative control is obtained by taking $b=0$, $q=0$, and
$c=e_0-e_3$. The three orbital sums of $c$ vanish, so
$\Pi_H(b,c,q)=0$, but the image relations fail.
The harmonic readout does not replace the compatibility check
for the channels.

\subsection{Conjugation, covariance, and memory of the incidence register}
\label{iii:r25:transporte}

Let $\rho(m)=13-m$. If incidence conjugation reverses the
order of the windows, preserves the arithmetic and boundary classes
of the visits, and interchanges the trace channels, then
\[
 (A_m^\dagger,C_m^\dagger,V_m^\dagger)
 =(A_{\rho(m)},-C_{\rho(m)},V_{\rho(m)}).
\]
With $c_m=[C_m]_{10}$, decimal normalization gives
\begin{equation}
 z_m^\dagger=z_{\rho(m)}
 +100\mathbf1_{\{c_{\rho(m)}\ne0\}}-20c_{\rho(m)}.
 \label{iii:r25:conjugacion}
\end{equation}
Indeed, $[-C]_{10}=0$ if $[C]_{10}=0$, and is
$10-[C]_{10}$ otherwise. The channel negation is
applied before reconstructing the decimal block.

Application to a specific involution also preserves the
window endpoints. For a route
$\gamma=(x_0,\ldots,x_n)$ read before advancement,
\begin{equation}
 \sum_{k=0}^{n-1}f(x_{n-k})
 =\sum_{k=0}^{n-1}f(x_k)+f(x_n)-f(x_0).
 \label{iii:r25:extremos}
\end{equation}
Both sums contain the same interior terms; the difference
is the final endpoint minus the initial one. On closed routes this
correction vanishes. On open windows it is incorporated before
reduction modulo $10$.

\begin{proposition}[Consequence of a periodic reduction]
\label{iii:r25:periodo}
If the observable of a fixed collection of routes satisfies
$o(t+54)=o(t)$, its multiplicities are preserved under that return,
and the same local readout acts on windows of nine times,
then
\[
 (A,C,V)_{m+6}=(A,C,V)_m,\qquad z_{m+6}=z_m.
\]
\end{proposition}
\begin{proof}
Translation $t\mapsto t+54$ takes each window bijectively
to the one six positions later and preserves each counted incidence
and its multiplicity. Equality is transmitted to the counts
and their decimal normalizations.
\end{proof}

A presentation with twelve distinct blocks therefore preserves
information that does not factor through that observable of period $54$:
incidence selection, orientation, endpoints, multiplicity, or
memory. The proposition delimits the effect of that reduction; it does
not assign period $54$ to the complete history.

\subsubsection{Determination on orbits and accumulation of events}

Fixing only $Q(z)=q$ leaves an affine class of the lattice
\[
 \mathcal Z_0=\{v\in\mathbb Z^{12}:\sum_i v_i=0\}
 =\bigoplus_{i=0}^{10}\mathbb Z(e_i-e_{11}).
\]
For example, $100\mathbf1+e_0$ and $100\mathbf1+e_1$ have
charge $1201$ and satisfy the image checks through their
respective channels. They are witnesses to the compatibility equations,
without any assertion that they belong to HMT terminal histories.

\begin{proposition}[Cyclic covariance]
\label{iii:r25:covariancia}
Let $\rho$ be the window translation of a register $R$, of
minimal period $p\mid12$, and let $S$ be the translation of its
coordinates. A covariant readout on that orbit is determined
by $z\in\operatorname{Fix}(S^p)$. The condition $Q(z)=q$
admits integer solutions if and only if $12/p$ divides $q$;
when they exist, they form an affine class of rank $p-1$.
\end{proposition}
\begin{proof}
The rule $F(\rho^jR)=S^jz$ is well defined exactly when
$S^pz=z$. These vectors consist of $12/p$ repetitions of a block of
length $p$, and their charge is $\frac{12}{p}\sum_{i=0}^{p-1}z_i$.
The formula proves the divisibility and rank. Covariance of
the channels follows from $D_aS=SD_a$ and $QS=Q$.
\end{proof}

In the harmonic chart, the corresponding action is
$\widehat S_H=\mathcal H_{12}S\mathcal H_{12}/4$.
If the origin is marked, the action also transports that mark.
The period of the orbit indicated here is not identified with that of
the enriched history. For an initial readout $F_{108}$ already
defined, naturality under truncation is expressed by
$F_N=F_{108}\circ\tau_{N,108}$; this identity preserves the
initial evaluation without selecting its production rule.

\begin{proposition}[Additive composition of incidence contributions]
\label{iii:r25:eventos}
Suppose that the initial state and each event $a_t$ produce,
through previously fixed rules on the register, contributions
$(\delta w_t,\delta q_t)$ such that
\[
 \sum_i\delta w_{t,i}=0,
 \qquad \delta q_t+T(\delta w_t)\equiv0\pmod {12}.
\]
With compatible initial conditions $(w^{(0)},q^{(0)})$,
updating by summation determines a unique register at each stage.
Its increment is
\[
 \delta z_{t,i}=
 \frac{\delta q_t+T(\delta w_t)}{12}-P_i(\delta w_t).
\]
The construction respects the concatenation of segments and preserves
the readout of a prefix under subsequent extensions.
\end{proposition}
\begin{proof}
The compatibility conditions are preserved by summation, and the
inverse \eqref{iii:r25:inversa} is additive in $(w,q)$.
Induction fixes each state. Summation over concatenated segments
is the sum of their contributions; a prefix receives the same
summands before and after extension.
\end{proof}

This proposition gives a sufficient criterion for additive
accumulation, not a requirement that every TPK readout be linear.
When events are transported by nontrivial actions, their
contributions are expressed in the terminal chart through the
corresponding affine cocycle. In both cases, the event rules and the
initial state preserve their incidence provenance.

The preceding development determines the passage
\[
 \mathscr L\longmapsto(A,C,V)\longmapsto z
 \longmapsto(D_3z,D_4z,Q(z))\longmapsto\mathcal H_{12}z
\]
with the conditions of each map. Its domain differs from the
regional register in Appendix~\ref{iii:nucleo-finito}: there, the
equations $X_aL_a=Y_a$ determine the lifts from
ternary transitions. Here, the incidences of visits and traces
determine integer blocks and their channels. The uniqueness of one
reconstruction does not replace the production of the register used
by the other. Both preserve their primitives, operations, and proofs
in the APP--TRIT--TPK chain.
